\section{Evaluation}
\label{sec:evaluation}

A machine-checked theorem can hold for uninteresting reasons. Its premises may
never hold together, its obligations may exclude no error a real analyzer
makes, and an analyzer that answers \verdict{UNKNOWN} everywhere satisfies it
as well (\isaname{unknown_everywhere_sound}). We therefore check that the
premises can hold and that the obligations have force, and then measure
precision on concrete programs. Facts about one fixed program are proved by
\texttt{eval}, which runs generated code and therefore also trusts the code
generator.

\subparagraph*{The theorem is not vacuous (machine-checked).}
For the program of \cref{fig:running}, analyzed with Interval, entry-state
contexts, and the warrowing update rule, a theory builds a source execution
that returns from both calls, evaluation shows that the analyzer answers with
both checks \verdict{PROVED}, and \isaname{nv_source_certified} instantiates
\cref{thm:headline} with every premise discharged. Companion instances do the
same for single checks (\isaname{nv_check_proved_sound}) and for a
\verdict{DEAD} check (\isaname{nv_dead_unreached}).

\subparagraph*{The obligations reject a real defect (machine-checked).}
In the Goblint defect of \cref{sec:intro}, the congruence domain returned the
constant~$1$ for $(1 + 2\mathbb{Z}) \bmod 2$~\cite{goblint1161}. In Voblint,
every congruence remainder must satisfy \isaname{congruence_mod_sound}, which
requires the truncating remainder of any two concrete values to lie in the
concretization of the abstract result. The pre-fix answer violates this at
$-5$, and \isaname{prefix_congruence_mod_unsound} proves that every abstract
remainder returning the constant does. The defective operator could not have
been registered as a domain. \oblig{Total} has force as well:
\isaname{total_dropped_unsound} exhibits a cover that meets the other four
obligations and misses a reached store.

\subparagraph*{Contexts buy precision (machine-checked for one program).}
When \texttt{f} passes its argument to \texttt{g} and \texttt{main} calls
\texttt{f} with a positive and a negative value, call strings of length one
merge both activations of \texttt{g} and Sign loses both signs, while length
two keeps them apart and proves both checks. The theories prove the separation
at the entry of \texttt{g}
(\isaname{sign_k2_strictly_more_precise_than_k1_at_g}). This is a witness for
one program, not a theorem that longer call strings dominate.

\subparagraph*{Comparison with Goblint (empirical).}
Programs chosen to show a mechanism succeeding favor Voblint, so we also use
programs selected by someone else. Monniaux's tutorial on the flavors of
abstract interpretation gives small programs on which standard analyses gain
or lose precision~\cite{monniaux24}. Every program VIMP can express was
transliterated into a regression fixture, and every check holds in every
execution, so each missing proof is lost precision. \Cref{tab:flavors}
compares Voblint's verdicts, which the build re-checks against the analyzer,
with Goblint at revision \texttt{6127869c} on C versions of the same programs,
each analyzer under the settings that give it its best chance.

\begin{table}[t]
\centering
\scriptsize
\setlength{\tabcolsep}{3.2pt}
% generated by paper/tools/flavors.py
\begin{tabular}{llcccccccc}
\toprule
 & & signs & not ind. & abs. & bounded & reset & reset & skipped & reset eq. \\
 & & \texttt{z < 0} & \texttt{i > 0} & \texttt{x != 0} & \texttt{i == 10} & \texttt{i <= 9} & \texttt{i >= 0} & \texttt{i <= 9} & \texttt{i <= 9} \\
\midrule
\multirow{4}{*}{Voblint} & Sign & \checkmark & -- & -- & -- & -- & \checkmark & -- & -- \\
 & Interval & \checkmark & -- & -- & \checkmark & \checkmark & \checkmark & -- & -- \\
 & Int & \checkmark & \checkmark & -- & \checkmark & \checkmark & \checkmark & -- & -- \\
 & Int, Order & \checkmark & \checkmark & -- & \checkmark & \checkmark & \checkmark & -- & -- \\
\midrule
\multirow{7}{*}{Goblint} & Interval & \checkmark & -- & -- & \checkmark & \checkmark & -- & -- & -- \\
 & delay 10 & \checkmark & -- & -- & \checkmark & \checkmark & \checkmark & \checkmark & \checkmark \\
 & Int, delay 3 & \checkmark & \checkmark & -- & \checkmark & \checkmark & \checkmark & -- & -- \\
 & thresholds & \checkmark & -- & -- & \checkmark & \checkmark & \checkmark & \checkmark & \checkmark \\
 & comparisons & \checkmark & -- & -- & \checkmark & \checkmark & \checkmark & -- & -- \\
 & octagons & \checkmark & -- & -- & \checkmark & \checkmark & -- & -- & -- \\
 & octagons, paths & \checkmark & -- & \checkmark & \checkmark & \checkmark & -- & -- & -- \\
\bottomrule
\end{tabular}

\caption{Checks proved (\checkmark) on the flavors programs~\cite{monniaux24}.
Voblint runs its domains, the reduced product Int, and Int with the order
analysis. Goblint's rows add a widening delay, widening thresholds, Interval
refined by congruences, and Apron octagons, the last also partitioned by the
branches taken. Each check holds in every execution.}
\label{tab:flavors}
\end{table}

Three results stand out. Int proves the non-inductive check \texttt{i > 0},
which neither Sign nor Interval proves alone: widening drops Interval's lower
bound, but Sign keeps $i \geq 0$, and the reduction recovers an even value in
$[2, 4]$. Goblint proves it only when Interval is refined by congruences, with
a widening delay of~3 or with thresholds at the constants of comparisons. The absolute
value example needs the two branches kept apart, which only Goblint's
partitioned octagons do; Voblint has no partitioning, which would need its own
soundness proof. The two variants on which narrowing is foiled (skipped
increment, reset at equality) need a widening delay or thresholds, which
Voblint's solver does not offer. Voblint's Interval proves \texttt{i >= 0} on
the reset example where Goblint's does not, because Goblint's widening stops at
the largest \texttt{int} and its default overflow handling then gives
\texttt{i + 1} the whole range; VIMP's integers are unbounded.

\subparagraph*{Running time (empirical).}
Speed was not a goal: the generated code runs as Isabelle emits it, with one
proved code equation added to remove a recompilation of the program at every
equation evaluation. Over the \stat{prof.fixtures} regression fixtures under
\stat{prof.configurations} combinations of domain and context policy, three
repetitions each on an Apple M4 Pro, the median solve takes
\stat{prof.solve_median_ms}\,ms and the 95th percentile
\stat{prof.solve_p95_ms}\,ms; \stat{prof.timeouts} of \stat{prof.runs} runs
exceed the 30\,s limit. The expensive cases include recursive programs under entry-state contexts,
where domains such as Interval create one context per recursion level.
On generated programs that grow one size at a time, the number of solver
evaluations grows about linearly, while solve time grows with exponents between
\stat{prof.scaling.time_min} and \stat{prof.scaling.time_max}, because the
generated code represents finite sets and maps as lists. These measurements
cover small programs on one machine and imply no performance bound.
